\addcontentsline{toc}{subsubsection}{Solution 52. The Almost-Cosine Sequence}

\begin{dlsummary}[title={\bfseries Solution 52. The Almost-Cosine Sequence},breakable]

\begin{proof}
The unique solution is
\[
\boxed{
f(n)=\frac{(2+\sqrt3)^n+(2-\sqrt3)^n}{2}
\qquad(n\in\mathbb Z).
}
\]

\paragraph{1. Force growth.}

Taking \(m=1,n=0\),
\[
|4-4f(0)|\le1.
\]
Since \(f(0)\) is an integer, \(f(0)=1\).

Now take \(m=n\):
\[
f(2n)\ge2f(n)^2-2.
\]
Starting from \(f(1)=2\), this gives
\[
f(2)\ge6,\qquad f(4)\ge70,
\]
and repeated doubling proves that \(f\) is unbounded.

\paragraph{2. Introduce the error.}

Write
\[
E(a,b)=f(a+b)+f(a-b)-2f(a)f(b),
\qquad |E(a,b)|\le1.
\]

Consider
\[
S=f(x+y+z)+f(x+y-z)+f(x-y+z)+f(x-y-z).
\]

Pairing terms with centres \(x+y\) and \(x-y\), we obtain
\[
S=
4f(x)f(y)f(z)
+2f(z)E(x,y)
+E(x+y,z)+E(x-y,z).
\]

Pairing instead with centre \(x\) and offsets \(y+z\) and \(y-z\), we obtain
\[
S=
4f(x)f(y)f(z)
+2f(x)E(y,z)
+E(x,y+z)+E(x,y-z).
\]

Subtracting these expressions gives
\[
2f(x)|E(y,z)|\le2f(z)+4.
\]
For fixed \(y,z\), the right side is constant, while \(f(x)\) is unbounded. Therefore
\[
E(y,z)=0.
\]
Thus the original inequality secretly forces
\[
\boxed{f(m+n)+f(m-n)=2f(m)f(n).} \tag{1}
\]

\paragraph{3. Recover the sequence.}

Taking \(m=0\) in (1), using \(f(0)=1\), gives
\[
f(-n)=f(n).
\]

Taking \(n=1\), using \(f(1)=2\), gives
\[
f(m+1)=4f(m)-f(m-1).
\]
Consequently, \(f(0)=1\) and \(f(1)=2\) determine every value:
\[
1,\ 2,\ 7,\ 26,\ 97,\ 362,\ldots
\]

The roots of \(r^2-4r+1=0\) are \(2+\sqrt3\) and \(2-\sqrt3\). Solving the recurrence gives the stated formula.

\paragraph{4. Verify.}

Put \(A=2+\sqrt3\), so \(A^{-1}=2-\sqrt3\). Direct expansion shows that
\[
f(n)=\frac{A^n+A^{-n}}2
\]
satisfies (1), hence the original inequality with zero error.

These values are positive, and the recurrence with integer initial values proves they are integers for nonnegative \(n\). Evenness handles negative \(n\).
\end{proof}
\end{dlsummary}

\begin{takeaways}
\begin{itemize}
\item \textbf{Find growth before controlling errors:} Unboundedness supplies the amplification needed to eliminate a fixed discrepancy.
\item \textbf{Regroup a common expression:} The four-term sum links errors at different input pairs.
\item \textbf{Separate rigidity from classification:} First establish the exact equation; only then solve its recurrence.
\item \textbf{Irrational formulas can give integers:} A recurrence often explains integrality more clearly than a closed form.
\end{itemize}
\end{takeaways}
